% ===== 以下为模板使用手册相关术语 =====
\newglossaryentry{cursor}{
    name={Cursor},
    first={Cursor},
    description={以 AI 辅助编程为核心的编辑器环境，可用于本模板的本地写作与编译管理。},
    sort={cursor}
}

\newglossaryentry{vscode}{
    name={VS Code},
    first={VS Code},
    description={常用代码编辑器，可配合 LaTeX 插件完成论文模板编辑与构建。},
    sort={vs code}
}

\newglossaryentry{overleaf}{
    name={Overleaf},
    first={Overleaf},
    description={在线 LaTeX 编辑平台，适合轻量协作与预览，但大型工程可能出现编译耗时较长。},
    sort={overleaf}
}

\newglossaryentry{texlive}{
    name={TeX Live},
    first={TeX Live},
    description={主流 TeX 发行版之一，提供论文模板编译所需工具链。},
    sort={tex live}
}

\newglossaryentry{miktex}{
    name={MiKTeX},
    first={MiKTeX},
    description={常用 TeX 发行版之一，可作为本模板本地编译环境。},
    sort={miktex}
}

\newglossaryentry{xelatex}{
    name={XeLaTeX},
    first={XeLaTeX},
    description={支持 Unicode 与系统字体的 LaTeX 引擎，适合中文论文排版。},
    sort={xelatex}
}

\newglossaryentry{latexworkshop}{
    name={LaTeX Workshop},
    first={LaTeX Workshop},
    description={VS Code/Cursor 常用 LaTeX 扩展，支持构建链、日志定位与 PDF 预览。},
    sort={latex workshop}
}

\newglossaryentry{bibtex-tool}{
    name={BibTeX},
    first={BibTeX},
    description={用于处理参考文献数据库并生成文献列表的工具。},
    sort={bibtex tool}
}

\newglossaryentry{tikz}{
    name={TikZ},
    first={TikZ},
    description={LaTeX 绘图库，可直接在文档中定义流程图、结构图等矢量图形。},
    sort={tikz}
}

\newglossaryentry{pgfplots}{
    name={PGFPlots},
    first={PGFPlots},
    description={基于 TikZ 的绘图包，常用于函数图和实验结果可视化。},
    sort={pgfplots}
}

\newglossaryentry{cleverefpkg}{
    name={cleveref},
    first={cleveref},
    description={交叉引用增强宏包，可自动补充图、表、公式和章节等引用前缀。},
    sort={cleveref package}
}

\newglossaryentry{glossariespkg}{
    name={glossaries},
    first={glossaries},
    description={术语表与缩略语管理宏包，可统一术语首次出现与索引输出。},
    sort={glossaries package}
}